
% !TEX program = pdflatex
\documentclass[11pt,a4paper]{article}
\usepackage[margin=2.5cm]{geometry}
\usepackage{times}
\usepackage{hyperref}

\title{Robust Deep Learning for Vision in Autonomous Robotics: Context, Literature, and Critical Discussion}
\author{}
\date{}

\begin{document}
\maketitle

\section{Introduction}
Deep neural networks (DNNs) have become the de facto standard for visual perception in autonomous robots, from self-driving cars to drones and manipulators, yet their brittleness to sensor noise, weather, and other real-world corruptions poses a fundamental safety risk. Unlike internet-scale photo tagging, where a misclassification merely produces an incorrect label, failures in robotic vision can directly cause collisions, unsafe manoeuvres, or mission aborts. This literature review synthesises recent work on robustness and sim-to-real transfer for vision-based robotic control, arguing that the distribution of training data---and particularly the design of augmentations and domain randomisation---often matters more for safe deployment than marginal improvements in clean-dataset accuracy or modest architectural tweaks.

Deep learning is now deeply embedded in robotic perception stacks, enabling tasks such as object detection, semantic segmentation, depth estimation, and end-to-end visuomotor control. In autonomous driving, deep detectors and segmenters are used to recognise vehicles, pedestrians, and traffic signs, while in manipulation and mobile robotics, end-to-end policies often map raw images directly to motor commands. This centrality of deep vision means that robustness properties of these models are no longer a purely academic concern: perception failures translate into physical risk rather than reduced benchmark scores.

A crucial contextual distinction is between vision for disembodied recognition (for example, ImageNet classification) and vision tightly coupled to control loops in safety-critical environments. In the former, common benchmarks evaluate accuracy on clean, well-curated datasets, whereas in the latter, the environment is dynamic, sensors are imperfect, and adverse weather or motion blur is routine rather than exceptional. This mismatch has led to a growing recognition that standard computer vision metrics and evaluation protocols are insufficient for robotics, where rare but catastrophic failures under distribution shift matter more than average-case performance.

\section{State of the Art and Challenges}

\subsection{The sim-to-real gap and structural brittleness}
Training robotic policies directly on real hardware is expensive, slow, and risky, so many state-of-the-art systems rely heavily on simulation for data generation and policy learning. The resulting ``sim-to-real gap'' arises because simulators rarely match the real world in visual appearance (textures, lighting, sensor artefacts) or dynamics (friction, delays, wear), leading to severe degradation when policies or perception models are deployed outside their training domain. Survey work on sim-to-real transfer in deep reinforcement learning formalises this problem and documents numerous failures where policies that perform well in simulation collapse in the real world once they encounter unmodelled disturbances or unseen visual configurations.

This gap is exacerbated by the structural brittleness of deep models to even modest corruptions. Benchmarking studies on object detection and autonomous driving show that standard CNN-based detectors lose a large fraction of their performance under synthetic corruptions such as noise, blur, fog, rain, and snow, despite remaining near state-of-the-art on clean images. For autonomous driving in particular, robustness evaluations demonstrate similar fragility when lidar or camera inputs are degraded by adverse weather and sensor corruptions, underscoring that current perception stacks are not inherently robust.

\subsection{How standard CNNs fail under noise}
A recurring theme across robustness studies is that conventional CNNs trained on clean data exhibit systematic, not random, failure patterns in the presence of corruptions. Classification and detection models trained on standard datasets experience a sharp drop in accuracy when evaluated on corrupted variants, with some corruptions---such as heavy snow, fog, and certain types of noise---producing disproportionate degradation. These performance drops are not merely small perturbations; they often correspond to models becoming effectively unusable at higher corruption severities, which in robotics translates directly to loss of situational awareness.

Application-focused studies in intelligent transportation systems reach similar conclusions. Robustness analyses of popular detectors such as YOLO and Faster R-CNN show that architectures which rank highly on clean benchmarks significantly underperform in noisy or low-visibility conditions compared to their clean baselines, emphasising that safety-critical perception requires explicit robustness-focused training and evaluation. Empirical work on training data perturbations further finds that robustness cannot be inferred from clean accuracy alone and that models optimised solely for clean performance may be especially vulnerable to random corruptions and adversarial perturbations.

Beyond empirical evidence, recent analyses of more advanced vision architectures, including transformer-based systems, show that even these models exhibit distinct vulnerability patterns tied to frequency biases and corruption types. This suggests that structural brittleness is a broad property of current deep vision models rather than a quirk of specific CNNs. Collectively, these findings support the view that robotic vision cannot rely on clean-benchmark accuracy as a proxy for safety; robustness to realistic noise, weather, and sensor imperfections must be explicitly engineered.

\subsection{Data augmentation and corruption-aware training}
One major response to brittleness has been to modify the training data distribution through targeted augmentations and corruption-aware training schemes. Classical techniques such as adding Gaussian noise and blur, as well as more sophisticated augmentation strategies like stylisation-based augmentation and model-based corruption generation, have been shown to improve robustness to a range of corruptions. For example, studies on corruption benchmarks demonstrate that augmenting training data with stylised or corrupted images significantly enhances robustness, though often at some cost to clean accuracy, highlighting the central trade-off between robustness and nominal performance.

More recent work introduces structured corruption models that generate realistic weather effects (snow, rain) and sensor artefacts, then uses them in training loops to improve robustness in traffic-sign recognition and related tasks. Empirical studies on training data perturbations further support the idea that controlled noise during training acts as an effective regulariser, improving both random and adversarial robustness while keeping clean-data performance largely intact when perturbations are well-calibrated. These findings align with broader observations in the robustness literature: exposing models to a carefully designed spectrum of perturbations during training fundamentally reshapes their decision boundaries, making them less sensitive to spurious high-frequency features and more reliant on stable cues.

However, corruption-aware training also has limitations. Many augmentation schemes focus on a fixed set of synthetic corruptions, and robustness gains do not always transfer to out-of-distribution distortions or to the full complexity of real sensor noise and weather conditions. In robotics, where the space of possible visual and physical disturbances is vast, this raises the question of how to design augmentations that meaningfully approximate deployment environments rather than simply overfitting to benchmark corruptions.

\subsection{Domain randomisation and visual sim-to-real transfer}
Domain randomisation has become one of the most influential strategies for bridging the sim-to-real gap in vision-based robotics. The core idea is to randomise visual properties---textures, lighting, camera parameters, and sometimes even object geometries---so extensively in simulation that the real world appears as just another sample from the training distribution. Early work on object localisation for robotic manipulation and quadrotor control showed that controllers trained with heavily randomised textures and appearances could be transferred directly to real-world settings without real-image pretraining, establishing a proof of concept for visual domain randomisation.

More recent work provides systematic explorations of visual domain randomisation for robotic manipulation, tuning randomisation parameters on proxy tasks and then training visuomotor policies in simulation with those optimised settings. These policies achieve high success rates across diverse manipulation tasks when deployed on real robots, and importantly, they often outperform policies trained on limited real data, suggesting that a well-designed synthetic distribution can rival or exceed small real datasets for robustness. Survey work on sim-to-real transfer in deep reinforcement learning for robotics similarly highlights visual domain randomisation as a central tool for tasks including localisation, detection, pose estimation, and segmentation.

From a theoretical standpoint, domain randomisation has been analysed within a reinforcement learning framework in which the simulator is modelled as a family of Markov decision processes with tunable parameters. This analysis provides bounds on the sim-to-real gap and shows that transfer can succeed under mild assumptions, especially when policies are allowed to be history-dependent, effectively using memory to infer latent environment parameters. This theory supports the empirical intuition that randomisation should cover a sufficiently rich set of visual and physical parameters; otherwise, sim-to-real performance remains fragile.

\section{Critical Conclusion of the Literature}
Taken together, the literature on corruption robustness, noisy training, domain randomisation, and sim-to-real transfer converges on a key narrative: for vision in autonomous robotics, the training data distribution is often more decisive for real-world robustness than modest variations in network topology. Corruption benchmarks for classification and detection show that standard architectures suffer dramatic performance loss under noise, blur, and weather, and that relatively simple changes to training data, such as stylised images or corruption-based augmentations, can yield large robustness gains without altering the network structure. Empirical studies of noisy training data and perturbation-based retraining demonstrate that exposing convolutional neural networks to controlled noise during training acts as a strong regulariser, enhancing robustness to both random and adversarial perturbations, again with little architectural change.

The success of domain randomisation in sim-to-real transfer for manipulation, quadrotor control, and other tasks indicates that the breadth and realism of the simulated visual distribution largely dictates whether a given architecture will survive deployment in the wild. Theoretical analyses of domain randomisation formalise conditions under which sim-to-real transfer works, making explicit that coverage and diversity in the training domain, coupled with memory in the policy, are central to closing the gap. Surveys and case studies in autonomous driving and sim-to-real reinforcement learning repeatedly emphasise that safety and robustness issues stem from distribution shift and unmodelled corruptions more than from intrinsic limitations of CNN or transformer architectures.

This is not to say that architecture is irrelevant---robustness-focused designs, such as models incorporating biological priors or spatio-temporal fusion, can offer additional gains. However, in many practical robotic systems, the dominant failure mode is that the training data fails to reflect deployment conditions rather than that the network lacks capacity or expressiveness. Consequently, research that focuses solely on achieving marginally higher clean accuracy with new architectures, without rethinking training distributions and evaluation protocols, risks optimising the wrong objective for safety-critical robotics.

For autonomous robots whose failures have physical consequences, evaluation and design must therefore pivot from optimising clean accuracy to prioritising structural robustness, with the training data distribution treated as a first-class design parameter rather than an afterthought. This shift in perspective aligns the objectives of deep learning research with the practical demands of safety-critical robotics, where the cost of brittle perception is measured not in benchmark points but in collisions, damage, and eroded trust.

\bibliographystyle{agsm}
\bibliography{references}

\end{document}
